\chapter{Questionário Pós-Sessão}
\label{apendice:questionario}

Este apêndice reproduz o questionário pós-sessão dos estudos empíricos no Brasil e em Angola. O instrumento contém 25 itens em escala Likert de 5 pontos (1 = Discordo Totalmente, 5 = Concordo Totalmente), organizados em quatro dimensões. No Brasil, responderam os dois braços do piloto: tutor ($n = 45$) e cartilha ($n = 12$). No braço da cartilha, os enunciados que mencionam o agente, o WhatsApp ou a gamificação (itens 6 a 13, 17, 20 e 24) foram adaptados à leitura, com as referências ao agente substituídas; a redação adaptada não foi preservada, e a versão reproduzida abaixo é a do braço do tutor. Por isso, as médias de satisfação dos dois braços provêm de instrumentos distintos e não são diretamente comparáveis. Em Angola, os indicadores de percepção referem-se apenas ao braço do tutor ($n = 22$; Seção~\ref{sec:q_pilotos}).

O Capítulo~\ref{cap:resultados_preliminares} usa as dimensões de motivação e de resolução de dúvidas de duas formas. No contraste brasileiro entre tutor e cartilha, cada participante é classificado como positivo quando a média individual da dimensão é igual ou superior a 4, e o indicador é a proporção de participantes positivos. Na comparação entre países, o indicador é a média da dimensão expressa em percentual do máximo (por exemplo, 90\% corresponde a 4,50/5), apresentada sem teste. A satisfação é a média da dimensão na escala de 1 a 5.

O questionário mede \emph{percepção}. Aplicado uma única vez, ao final da sessão, ele não afere conhecimento financeiro, e mesmo os itens da Dimensão~4 registram apenas a aprendizagem declarada pelo participante. O ganho de aprendizagem do Capítulo~\ref{cap:resultados_preliminares} provém do teste do Apêndice~\ref{apendice:teste_conhecimento}.

\section*{Dimensão 1: Motivação e Engajamento}

\begin{enumerate}
    \item Senti-me motivado(a) a participar das atividades de educação financeira propostas pelo agente.
    \item A interação com o agente me incentivou a buscar mais informações sobre finanças pessoais.
    \item Senti-me engajado(a) durante toda a sessão de interação.
    \item As atividades propostas despertaram meu interesse em melhorar minha situação financeira.
    \item Gostaria de continuar utilizando o sistema para aprender mais sobre educação financeira.
    \item A gamificação (sistema de pontos) me motivou a completar as missões.
\end{enumerate}

\section*{Dimensão 2: Satisfação com a Interação}

\begin{enumerate}
\setcounter{enumi}{6}
    \item A experiência de conversar com o agente foi agradável.
    \item As respostas do agente foram claras e fáceis de entender.
    \item O agente respondeu de forma adequada às minhas perguntas e dúvidas.
    \item Senti que o agente compreendeu minhas necessidades financeiras.
    \item A duração da sessão foi adequada.
    \item Recomendaria este sistema a amigos ou familiares.
    \item A utilização do WhatsApp facilitou a interação com o sistema.
\end{enumerate}

\section*{Dimensão 3: Adequação Cultural Percebida}

\begin{enumerate}
\setcounter{enumi}{13}
    \item A linguagem utilizada pelo agente foi natural e adequada ao meu contexto.
    \item Os exemplos financeiros apresentados foram relevantes para a minha realidade.
    \item Os instrumentos financeiros mencionados (tipos de conta, formas de pagamento) são os que eu utilizo no dia a dia.
    \item O agente demonstrou compreensão sobre a realidade econômica do meu país.
    \item Senti que o conteúdo foi adaptado para pessoas na minha situação financeira.
    \item Os valores monetários e cenários apresentados foram realistas para o meu contexto.
\end{enumerate}

\section*{Dimensão 4: Resolução de Dúvidas Financeiras}

\begin{enumerate}
\setcounter{enumi}{19}
    \item O agente me ajudou a compreender melhor minha situação financeira atual.
    \item Aprendi conceitos financeiros novos durante a interação.
    \item As ferramentas apresentadas (orçamento, plano de ação) são úteis para minha vida.
    \item Sinto-me mais preparado(a) para tomar decisões financeiras após a interação.
    \item O agente conseguiu resolver todas as minhas dúvidas sobre finanças pessoais.
    \item Pretendo aplicar na prática o que aprendi durante a sessão.
\end{enumerate}

\vspace{1cm}

\noindent\textbf{Escala de resposta:}\\
1 = Discordo Totalmente \quad 2 = Discordo \quad 3 = Neutro \quad 4 = Concordo \quad 5 = Concordo Totalmente

%%%%%%%%%%%%%%%%%%%%%%%%%%%%%%%%%%%%%%%%%%%%%%%%%%%%%%%%%%%%%%%%%%%%%%%%%%%%%%%%

\chapter{Teste de Conhecimento Financeiro (Pré e Pós-Sessão)}
\label{apendice:teste_conhecimento}

Este apêndice documenta o instrumento de aferição de conhecimento financeiro aplicado antes e
depois da sessão nos dois estudos-piloto (Seção~\ref{sec:q_pilotos}): no Brasil, aos braços do
tutor ($n = 45$) e da cartilha ($n = 12$); em Angola, aos braços do tutor ($n = 22$) e da
cartilha ($n = 8$). Os itens reproduzidos abaixo são os da versão brasileira; a versão
angolana tem os mesmos itens, com os valores expressos em kwanzas (Kz). Além da moeda, os
itens 5 (``pedidos de comida por aplicativo'') e 6 (``bicos'') carregam contexto brasileiro.
Não houve verificação de invariância entre as duas versões, e o ganho de Hake angolano não é
comparável ao brasileiro. Diferentemente do questionário do Apêndice~\ref{apendice:questionario}, que mede
\emph{percepção}, este instrumento mede \emph{desempenho}: cada item tem resposta correta única,
e o escore é a proporção de acertos.

O conteúdo cobre os três eixos da Missão~M1, Raio-X das Finanças
(Apêndice~\ref{apendice:missoes}), tarefa atribuída aos participantes: identificação de custos
fixos e gastos variáveis, apuração de receitas e leitura do resultado orçamentário. Os itens
distribuem-se entre \textbf{conhecimento declarativo} (reconhecer e definir conceitos) e
\textbf{conhecimento procedimental} (executar o cálculo ou aplicar a regra a um caso concreto),
distinção adotada para evitar que o escore refletisse apenas memorização terminológica. A
Tabela~\ref{tab:matriz_teste} associa cada item a um domínio, a um tipo de conhecimento e a uma
fase da M1.

\section*{Matriz de especificação}

\begin{table}[htbp]
\centering
\caption[Matriz de especificação do teste de conhecimento]{Matriz de especificação do teste de conhecimento por domínio, tipo de conhecimento e fase da Missão~M1.}
\label{tab:matriz_teste}
\begin{tabular}{clcc}
\hline
\textbf{Item} & \textbf{Domínio} & \textbf{Tipo de conhecimento} & \textbf{Fase da M1} \\
\hline
1  & Orçamento pessoal        & Declarativo    & Introdução ao Orçamento \\
2  & Custos fixos             & Declarativo    & Listar Custos Fixos \\
3  & Custos fixos             & Procedimental  & Listar Custos Fixos \\
4  & Gastos variáveis         & Declarativo    & Listar Gastos Variáveis \\
5  & Gastos variáveis         & Procedimental  & Listar Gastos Variáveis \\
6  & Receitas                 & Declarativo    & Informar Receitas \\
7  & Fluxo de caixa           & Declarativo    & Reflexão sobre o Resultado \\
8  & Fluxo de caixa           & Procedimental  & Reflexão sobre o Resultado \\
9  & Capacidade de poupança   & Procedimental  & Reflexão sobre o Resultado \\
10 & Planejamento financeiro  & Procedimental  & Reflexão sobre o Resultado \\
\hline
\end{tabular}

{\footnotesize Fonte: elaborado pelo autor.}
\end{table}

\section*{Itens}

\noindent\textbf{1.} O que é um orçamento pessoal?
\begin{enumerate}
    \item[(a)] Um empréstimo que o banco oferece todo mês.
    \item[(b)] Um registro do quanto se recebe e do quanto se gasta em um período. \textbf{(correta)}
    \item[(c)] O valor total que uma pessoa tem guardado no banco.
    \item[(d)] Uma lista de dívidas em atraso.
\end{enumerate}

\vspace{0.3cm}
\noindent\textbf{2.} Qual das despesas abaixo é um \textbf{custo fixo}?
\begin{enumerate}
    \item[(a)] Lanche na padaria.
    \item[(b)] Aluguel da casa. \textbf{(correta)}
    \item[(c)] Presente de aniversário.
    \item[(d)] Conserto inesperado do carro.
\end{enumerate}

\vspace{0.3cm}
\noindent\textbf{3.} Dona Maria paga todo mês: aluguel de R\$~600, energia de R\$~120,
internet de R\$~80 e gastou R\$~50 num lanche no domingo. Qual é o total dos \textbf{custos
fixos} dela no mês?
\begin{enumerate}
    \item[(a)] R\$~600
    \item[(b)] R\$~800 \textbf{(correta)}
    \item[(c)] R\$~850
    \item[(d)] R\$~250
\end{enumerate}

\vspace{0.3cm}
\noindent\textbf{4.} O que caracteriza um \textbf{gasto variável}?
\begin{enumerate}
    \item[(a)] Muda de valor ou de frequência conforme o mês. \textbf{(correta)}
    \item[(b)] Tem sempre o mesmo valor todo mês.
    \item[(c)] É sempre um gasto desnecessário.
    \item[(d)] É cobrado apenas uma vez por ano.
\end{enumerate}

\vspace{0.3cm}
\noindent\textbf{5.} Seu Antônio quer reduzir gastos sem perder moradia nem transporte para o
trabalho. Qual despesa ele deve revisar \textbf{primeiro}?
\begin{enumerate}
    \item[(a)] O aluguel.
    \item[(b)] A passagem de ônibus para o trabalho.
    \item[(c)] Os pedidos de comida por aplicativo. \textbf{(correta)}
    \item[(d)] A conta de água.
\end{enumerate}

\vspace{0.3cm}
\noindent\textbf{6.} Uma pessoa recebe salário e ainda faz bicos aos fins de semana. Ao montar o
orçamento, ela deve considerar como \textbf{receita}:
\begin{enumerate}
    \item[(a)] Apenas o salário, porque é o valor garantido.
    \item[(b)] Apenas os bicos, porque é dinheiro extra.
    \item[(c)] O salário e os bicos, porque ambos entram no orçamento. \textbf{(correta)}
    \item[(d)] Nenhum dos dois, porque receita é só o que sobra no fim do mês.
\end{enumerate}

\vspace{0.3cm}
\noindent\textbf{7.} O que o \textbf{fluxo de caixa} mostra?
\begin{enumerate}
    \item[(a)] Somente o dinheiro que entrou.
    \item[(b)] Somente o dinheiro que saiu.
    \item[(c)] A diferença entre o que entrou e o que saiu no período. \textbf{(correta)}
    \item[(d)] O valor das dívidas acumuladas ao longo dos anos.
\end{enumerate}

\vspace{0.3cm}
\noindent\textbf{8.} No mês passado, a família da Joana recebeu R\$~2.400 e gastou R\$~2.700.
O fluxo de caixa dela foi:
\begin{enumerate}
    \item[(a)] Positivo em R\$~300.
    \item[(b)] Negativo em R\$~300. \textbf{(correta)}
    \item[(c)] Equilibrado.
    \item[(d)] Não é possível saber sem conhecer as dívidas.
\end{enumerate}

\vspace{0.3cm}
\noindent\textbf{9.} Carlos recebe R\$~1.800 por mês. Seus custos fixos somam R\$~1.100 e seus
gastos variáveis, R\$~400. Quanto ele consegue poupar por mês, mantendo os mesmos gastos?
\begin{enumerate}
    \item[(a)] R\$~700
    \item[(b)] R\$~400
    \item[(c)] R\$~300 \textbf{(correta)}
    \item[(d)] Nada, porque os gastos superam a renda.
\end{enumerate}

\vspace{0.3cm}
\noindent\textbf{10.} Uma pessoa descobre que gasta mais do que recebe todos os meses. Qual é a
\textbf{primeira} providência do planejamento financeiro?
\begin{enumerate}
    \item[(a)] Fazer um novo empréstimo para cobrir a diferença.
    \item[(b)] Levantar todas as receitas e despesas para saber onde está o desequilíbrio. \textbf{(correta)}
    \item[(c)] Parar de pagar as contas menos importantes.
    \item[(d)] Esperar o próximo aumento de salário.
\end{enumerate}

\section*{Aplicação e pontuação}

O instrumento foi aplicado duas vezes ao mesmo participante, imediatamente antes e
imediatamente após a sessão, na mesma forma e com a mesma ordem de itens. Cada item vale um
ponto; o escore do participante é a proporção de acertos, expressa em porcentagem. O ganho de
aprendizagem foi calculado pelo ganho normalizado de Hake~\cite{hake1998interactive},
$\langle g \rangle = (\text{pós} - \text{pré}) / (100 - \text{pré})$, que corrige o efeito de
teto para participantes com escore inicial mais alto.

\section*{Propriedades psicométricas e limitações}

Dois especialistas em educação financeira revisaram a correspondência entre cada item e o
domínio pretendido (validade de conteúdo), sem índice de concordância (CVI). A consistência
interna dos dez itens dicotômicos, medida pelo coeficiente KR-20 (equivalente ao $\alpha$ de
Cronbach para itens dicotômicos), foi de 0,78. A análise de itens (dificuldade e discriminação)
não foi feita. Três limitações condicionam a leitura: o mesmo instrumento nas duas aplicações
admite \emph{efeito de testagem}; os itens cobrem só os domínios da M1, e não a literacia
financeira em sentido amplo; e dez itens dicotômicos discriminam pouco entre escores próximos.
A análise de itens e as formas paralelas estão previstas na atividade A12
(Seção~\ref{sec:cronograma}).

\chapter{Estrutura das Missões Educacionais}
\label{apendice:missoes}

Este apêndice lista as fases das seis missões do sistema \textit{Jornada da Tranquilidade}: M1, Raio-X das Finanças; M2, Compra Inteligente; M3, Plano de Ação (as três obrigatórias); M4, Plano de Emergência; M5, Renegociação de Dívidas (as duas condicionais); e M6, Pesquisa de Preços (complementar). A pontuação por missão e as condições de ativação constam da Seção~\ref{sec:missoes}. Na Tabela~\ref{tab:missoes_completa}, as missões assinaladas com $\dagger$ (M1, M2, M3 e M6) têm fases que integram o espaço de ações $\mathcal{A}$ do agente de DRL, num total de 17 tarefas na especificação. Os experimentos do Capítulo~\ref{cap:resultados_preliminares} usaram um catálogo de 11 tarefas: 5 da M1, 2 da M2, 2 da M3 e 2 da M6 (Apêndice~\ref{apendice:experimentos}). As 8 fases de M4 e M5 ficam fora desse espaço e são geridas por regras condicionais de ativação.

\begin{table}[htbp]
\centering
\small
\caption[Fases das missões do sistema]{Fases das missões, na ordem de execução (A: ação, 75 pontos; R: reflexão, 25 pontos; máximo de 1.275 pontos).}
\label{tab:missoes_completa}
\begin{tabular}{lp{0.8\textwidth}}
\hline
\textbf{Missão} & \textbf{Fases} \\
\hline
M1$^{\dagger}$ & Introdução ao Orçamento (R); Listar Custos Fixos (A); Listar Gastos Variáveis (A); Informar Receitas (A); Reflexão sobre o Resultado (R) \\
M2$^{\dagger}$ & Necessidade vs.\ Desejo (R); Poupar vs.\ Parcelar (R) \\
M3$^{\dagger}$ & Introdução ao Plano de Ação (R); Definir O QUÊ e POR QUÊ (A); Definir QUEM, QUANDO, QUANTO (A); Definir COMO (A); Reflexão sobre o Plano (R) \\
M4 & Mapear Todas as Dívidas (A); Estratégia de Unificação (R); Ações Recomendadas (R) \\
M5 & Introdução à Renegociação (R); Mapear Credores e Dívidas (A); Contato e Negociação (A); Comprovante da 1\textordfeminine{} Parcela (A); Reflexão sobre o Compromisso (R) \\
M6$^{\dagger}$ & Introdução à Pesquisa de Preços (R); Pesquisar Preço (A); Compartilhar com a Comunidade (A); Buscar Preços da Comunidade (A); Reflexão sobre Economia (R) \\
\hline
\end{tabular}

{\footnotesize Fonte: elaborado pelo autor. Nota: $^{\dagger}$missão cujas fases integram o espaço de ações $\mathcal{A}$ do agente de DRL.}
\end{table}

%%%%%%%%%%%%%%%%%%%%%%%%%%%%%%%%%%%%%%%%%%%%%%%%%%%%%%%%%%%%%%%%%%%%%%%%%%%%%%%%

\chapter{Variantes Culturais das Missões Educacionais}
\label{apendice:variantes}

A classificação compara as variantes angolanas e portuguesas ao desenho brasileiro, segundo as oito dimensões de Bernal et al.~\cite{bernal1995} e a distinção entre superfície e estrutura profunda de Resnicow et al.~\cite{resnicow1999}. A Tabela~\ref{tab:variantes_veredito} reúne os doze vereditos. A comparação completa da Missão~1 ilustra os dois níveis; o detalhamento das demais missões está no repositório de validação dos artefatos, \url{https://github.com/paulocob/validationPaperMultiAgentFinancialLiteracy} (Seção~\ref{subsec:disponibilidade}).

Nos pilotos, a missão atribuída foi a M1, aplicada ao braço do tutor no Brasil ($n = 45$) e em Angola ($n = 22$; Seção~\ref{sec:q_pilotos}). Alguns participantes avançaram para outras missões; o excerto angolano do apêndice de exemplos de interações da tese refere-se à M6. A documentação de aplicação por missão em Angola é parcial. Nenhuma variante portuguesa foi aplicada a aprendizes: a cobertura de Portugal é exclusivamente por \textit{corpus}.

Cinco variantes, todas angolanas, são classificadas como profundas, e as sete restantes como de superfície. Duas classificações têm ressalva: a M4 angolana recebe veredito de desenho, porque não há base observacional verificada sobre a composição do endividamento, e a M5 portuguesa é de superfície incompleta, pois o regime local de renegociação ainda precisa ser incorporado. Como nenhuma variante foi comparada experimentalmente à versão não adaptada da mesma missão, a classificação descreve o desenho e não demonstra eficácia.

\begin{table}[htbp]
\centering
\small
\caption[Vereditos de adaptação por missão]{Dimensões do Modelo de Validade Ecológica alteradas em cada variante e nível de adaptação.}
\label{tab:variantes_veredito}
\setlength{\tabcolsep}{4pt}
\begin{tabular}{lllll}
\hline
\textbf{Missão} & \textbf{Angola: dimensões} & \textbf{Nível (AO)} & \textbf{Portugal: dimensões} & \textbf{Nível (PT)} \\
\hline
M1 & L, P, M, Cd, Cc, O, Mt, Cx & Profunda & L, Cd, Mt & Superfície \\
M2 & L, P, M, Cd, Cc, O, Mt & Profunda & L, M, Cd & Superfície \\
M3 & L, Cd, Cc, O, Mt, Cx & Profunda & L, Cd, Mt & Superfície \\
M4 & L, Cd, Cc, O, Mt & Profunda$^{*}$ & L, Cd & Superfície \\
M5 & L, Cd, O, Mt, Cx & Profunda & L, Cd & Superfície$^{\dagger}$ \\
M6 & L, Cd, Mt, Cx & Superfície & L, Cd & Superfície \\
\hline
\end{tabular}

\begin{minipage}{0.92\textwidth}
{\footnotesize Fonte: elaborado pelo autor. Nota: dimensões conforme Bernal, Bonilla e Bellido~\cite{bernal1995}: L, linguagem; P, pessoas; M, metáforas; Cd, conteúdo; Cc, conceitos; O, objetivos; Mt, métodos; Cx, contexto. Níveis conforme Resnicow et al.~\cite{resnicow1999}. $^{*}$Veredito de desenho e $^{\dagger}$superfície incompleta: ressalvas no texto acima.}
\end{minipage}
\end{table}

\section*{M1: Raio-X das Finanças}

\noindent\textbf{Objetivo invariante.} Construir o primeiro orçamento completo do aprendiz (receitas, despesas fixas e despesas variáveis) e produzir o saldo que condiciona a ativação das missões M4 e M5. As cinco fases da missão (Apêndice~\ref{apendice:missoes}) permanecem as mesmas nas três variantes; o que muda é a unidade de apuração e o conjunto de rubricas.

\begin{enumerate}
\item[(a)] \textbf{Brasil (referência; aplicada).} Instrumentos e termos: salário, ``bico'', conta de luz e de água, internet, fatura do cartão, parcela de financiamento. A apuração é mensal, e a informalidade brasileira, de 37,4\% dos ocupados no 2º trimestre de 2026~\cite{ibge2026pnad}, é tratada como renda complementar sobre uma base salarial. Enunciado ao aprendiz, tal como registrado no piloto (apêndice de exemplos de interações da tese): ``qual é a sua renda mensal? Pode incluir salário, bicos e qualquer outra fonte''.

\item[(b)] \textbf{Angola (aplicada).} Instrumentos e termos: \textit{zunga} e comércio de mercado, \textit{candongueiro} (transporte \textit{colectivo} informal e, por extensão, o seu condutor~\cite{lopes2006candongueiros}), contribuição à \textit{kixikila} e ao Instituto Nacional de Segurança Social. A variante introduz quatro mudanças. (i)~A apuração cobre os últimos sete dias e o número de dias com venda, porque a informalidade atinge 79,8\% dos ocupados (83,9\% em Benguela e 59,5\% em Luanda) e 62,4\% dos empregados têm contrato verbal~\cite{ineangola2024emprego}; a ``renda mensal'' pressupõe uma regularidade que a maioria dos aprendizes não tem. (ii)~A contribuição à \textit{kixikila}, associação rotativa de poupança e crédito~\cite{carlos2014kixikila,ardener1964}, é ao mesmo tempo saída de caixa e acumulação, e por isso não é registrada como despesa. (iii)~A \textit{deslocação} diária em \textit{candongueiro} é tratada como despesa fixa. (iv)~A fase de reflexão incorpora a inflação (variação homóloga de 12,42\% em março de 2026~\cite{ineangola2026ipcn}), que dá custo à guarda de kwanza em espécie. A contribuição à segurança social integra a missão porque o regime especial angolano (1,5\% mais 4\%, contra 3\% e 8\% no regime geral) abrange os vendedores ambulantes e de mercado, isto é, as próprias \textit{zungueiras}~\cite{maptss2026inss}.

\noindent\textit{Enunciado ao aprendiz:} ``Vamos fazer o raio-X do dinheiro da senhora. Em vez de começar pelo mês inteiro, vamos pelos últimos sete dias: em quantos desses dias houve venda, e quanto entrou em cada um? Depois somamos o que sai todos os dias --- o candongueiro para ir e voltar do mercado, a comida --- e só então o que sai uma vez por mês. E há uma pergunta que não é despesa nem sobra: quanto a senhora entrega à kixikila neste ciclo, e quando é a sua vez de receber?''

\item[(c)] \textbf{Portugal (proposta; não aplicada).} Instrumentos e termos: vencimento, subsídio de férias e subsídio de Natal, previstos nos artigos 263º e 264º do Código do Trabalho e imperativos, de que resultam catorze meses de salário por ano~\cite{ct2009}; prestação do crédito à habitação; MB WAY~\cite{sibs_mbway}; IRS. A mudança substantiva está na base anual da renda, de catorze meses em vez de doze, e por isso a fase de reflexão pergunta se o orçamento mensal repartiu pelos doze meses a renda anual de catorze salários, questão inexistente no desenho brasileiro. A variante pergunta também se os subsídios são recebidos diluídos ao longo do ano ou em dois montantes. O regime legal transitório de duodécimos (2013--2017) não foi prorrogado para 2018, e desde então, no setor privado, o pagamento em duodécimos depende de acordo entre empregador e trabalhador; a variante trata a opção como circunstância declarada pelo aprendiz.

\noindent\textit{Enunciado ao aprendiz:} ``Vamos fazer o raio-X das suas contas. Primeiro, uma pergunta que muda todo o resto: recebe os subsídios de férias e de Natal diluídos ao longo do ano ou em dois montantes? Depois disso, listamos as despesas fixas --- e começo pela prestação da casa, se tiver.''
\end{enumerate}

\noindent\textbf{Classificação.} A variante \textbf{angolana é adaptação de estrutura profunda}: altera conceitos, porque a unidade de renda é o fluxo diário de caixa, e não o salário mensal; objetivos, porque a meta é cobrir os dias sem venda, em vez de gerar sobra no mês; métodos, porque o levantamento é retrospectivo de sete dias, em lugar de uma tabela mensal de rubricas; e contexto, pela informalidade mais que dobrada e pela inflação de dois dígitos. O desenho ocidental de orçamento pressupõe horizonte de planejamento longo e renda previsível, o que não se verifica nesse público; trata-se da dimensão de percepção temporal que Parrish e Linder-VanBerschot~\cite{parrish2010} identificam como crítica em instrução multicultural. A simplificação do levantamento encontra apoio, ademais, no achado de Drexler, Fischer e Schoar~\cite{drexler2014} de que o treinamento por regras de bolso supera o treinamento contábil padrão entre participantes com menor escolaridade ou práticas iniciais mais frágeis; esse achado apoia a adaptação de \emph{método} ao público, e não a adaptação cultural em sentido étnico. A variante \textbf{portuguesa é adaptação de superfície}: substitui léxico e instrumentos e introduz uma correção de cálculo (renda anual de catorze salários, repartida por doze meses), mas mantém o sujeito individual, a renda formal previsível, a unidade mensal e o objetivo de sobra.

\section*{Uso da classificação pelo módulo de DRL}
\label{sec:variantes_fcult}

A função de adequação cultural $f_{\text{cult}}$ da recompensa (Equação~\ref{eq:f_cult} e Tabela~\ref{tab:fcult} do Capítulo~\ref{cap:proposta}) deriva da Tabela~\ref{tab:variantes_veredito} sem parâmetro novo: cada par (país, missão) recebe 1,0 quando a variante é de estrutura profunda e 0,5 quando é de superfície, e o Brasil, contexto de referência, recebe 1,0 em todas as missões. Resultam Angola com 1,0 em M1, M2 e M3 e 0,5 em M6, e Portugal com 0,5 em todas. M4 e M5 não integram o espaço de ações e não têm valor na função; por isso, as ressalvas sobre a M4 angolana e a M5 portuguesa não afetam $f_{\text{cult}}$. A função mede a \emph{disponibilidade} de conteúdo adaptado, e não a sua eficácia junto ao aprendiz. O seu teste consta da Seção~\ref{sec:q_cultural}.

%%%%%%%%%%%%%%%%%%%%%%%%%%%%%%%%%%%%%%%%%%%%%%%%%%%%%%%%%%%%%%%%%%%%%%%%%%%%%%%%

%%%%%%%%%%%%%%%%%%%%%%%%%%%%%%%%%%%%%%%%%%%%%%%%%%%%%%%%%%%%%%%%%%%%%%%%%%%%%%%%

\chapter{Hiperparâmetros e Configurações dos Experimentos de DRL}
\label{apendice:experimentos}

As configurações reunidas aqui correspondem aos experimentos do módulo de DRL cujos resultados constam da Seção~\ref{sec:q_drl}. Os treinamentos sobre o conjunto sintético original operam sobre as 2.000 transições caracterizadas a seguir e são reproduzidos por um único comando (seção Reprodutibilidade); o simulador cognitivo, usado como segundo ambiente de avaliação, é descrito na última seção. Uma nota histórica registra os defeitos da primeira execução.

\section*{Características do conjunto sintético}

O programa que gerou o conjunto não foi preservado. As propriedades abaixo foram apuradas diretamente nos dados.

\begin{itemize}
    \item \textbf{Perfis}: 25 perfis simulados (9 iniciantes, 11 intermediários e 5 avançados), com literacia entre 0,22 e 0,84 e idade entre 21 e 55 anos. O sentimento inicial é neutro em 10 perfis, motivado em 8, ansioso em 5 e cansado em 2. Cada perfil tem de 75 a 86 transições (mediana de 81), num total de 2.000. O perfil não é constante por identificador em todo o conjunto: \texttt{u\_001} aparece como intermediário em uma transição.
    \item \textbf{Ações}: catálogo de 11 tarefas. M1: introdução, custos fixos, gastos variáveis, informar renda e reflexão sobre o resultado; M2: introdução e decisão; M3: introdução e definição do objetivo; M6: introdução e prática de pesquisa de preços. A última tarefa da M6 é a única terminal. Por missão, as ações distribuem-se em M6, 634; M1, 619; M3, 423; e M2, 324. Em 64,8\% das transições, a ação é a tarefa seguinte na sequência linear do catálogo. Em 11,5\% dos estados, a tarefa corrente é terminal e a regra sequencial não tem predição.
    \item \textbf{Recompensa do gerador}: 90 valores distintos no intervalo $[0{,}05;\ 0{,}95]$, com média de 0,532. A média é de 0,700 quando a ação segue a sequência linear e de 0,223 quando não segue. Ela varia também com o sentimento do estado: 0,746 (motivado), 0,559 (neutro), 0,348 (ansioso), 0,252 (cansado) e 0,267 (frustrado).
    \item \textbf{Transição}: a tarefa corrente do estado seguinte é sempre a ação tomada (2.000 de 2.000). O registro do estado seguinte não traz o perfil nem as missões concluídas, e 88 estados têm o campo de missões concluídas inconsistente com as tarefas concluídas. O codificador corrigido (seção seguinte) deriva os dois campos das tarefas concluídas. Com essa correção, a recompensa refinada tem $r_{\text{miss}} = 1$ em 3,1\% das transições e média de 0,516; o episódio termina em 8,3\% delas. Apenas 1.026 dos 2.000 estados seguintes reaparecem como estado do mesmo perfil, de modo que o conjunto não forma trajetórias encadeadas.
    \item \textbf{Repetição de estados e vazamento}: há 886 estados distintos em 2.000. Na partição original, por interação, 66,2\% dos estados de teste e 61,0\% dos pares $(\mathbf{s}, a)$ de teste aparecem idênticos no treino; na partição por perfil, 1,6\% e 1,6\%.
    \item \textbf{Justificativa da ação}: o campo de justificativa é texto livre, com 328 valores distintos, e não entra no estado.
\end{itemize}

\section*{Codificador de estado, partições por perfil e avaliação \textit{off-policy}}
\label{sec:apendice_particoes}

\noindent\textbf{Codificador corrigido.} Como o registro do estado seguinte não traz o perfil nem as missões concluídas, o codificador atribui a $\mathbf{s}_{t+1}$ o perfil da transição e deriva as missões concluídas das tarefas concluídas, tanto em $\mathbf{s}_t$ quanto em $\mathbf{s}_{t+1}$. Desse modo, as médias por grupo de atributos de $\mathbf{s}_{t+1}$ coincidem com as de $\mathbf{s}_t$ nos atributos de perfil (literacia 0,547; segmento e sentimento idênticos) e refletem o progresso nos demais. A nota histórica descreve o defeito que essa correção elimina.

\noindent\textbf{Partições por perfil.} Os 25 perfis foram divididos em cinco grupos disjuntos, estratificados por segmento, por rodízio após embaralhamento com semente 42:
\begin{itemize}
    \item grupo 0: \texttt{u\_002}, \texttt{u\_004}, \texttt{u\_006}, \texttt{u\_007}, \texttt{u\_008};
    \item grupo 1: \texttt{u\_012}, \texttt{u\_014}, \texttt{u\_016}, \texttt{u\_017}, \texttt{u\_020};
    \item grupo 2: \texttt{u\_001}, \texttt{u\_003}, \texttt{u\_015}, \texttt{u\_018}, \texttt{u\_024};
    \item grupo 3: \texttt{u\_005}, \texttt{u\_010}, \texttt{u\_011}, \texttt{u\_013}, \texttt{u\_023};
    \item grupo 4: \texttt{u\_009}, \texttt{u\_019}, \texttt{u\_021}, \texttt{u\_022}, \texttt{u\_025}.
\end{itemize}
No \emph{split} fixo, o teste é o grupo 0, a validação é o grupo 1 e o treino reúne os grupos 2 a 4 (5, 5 e 15 perfis; 375, 418 e 1.207 transições). Na validação cruzada de cinco \emph{folds}, o \emph{fold} $k$ usa o grupo $k$ como teste, o grupo $k+1$ como validação (parada antecipada e seleção de $\lambda_{\text{reg}}$, de forma aninhada) e os três grupos restantes como treino. O \emph{split} fixo foi executado com as cinco sementes 42 a 46, para medir a estabilidade entre inicializações; a validação cruzada, com a semente 42, e é ela que fornece a tabela principal (Tabela~\ref{tab:q_final}).

\noindent\textbf{Regra sequencial.} A política de referência é determinística: recomenda a tarefa seguinte à tarefa corrente no catálogo. Nos 11,5\% de estados cuja tarefa corrente é terminal, ela não tem predição, e o estado é contado como erro. A variante que repete a tarefa terminal nesses estados é reportada à parte (Top-1 de $0{,}776 \pm 0{,}031$ e recompensa associada de 0,627 na validação cruzada).

\noindent\textbf{Avaliação \textit{off-policy}.} Duas estimativas de valor acompanham as métricas de fidelidade. A importância amostral auto-normalizada (SNIS) usa a política de comportamento $\pi_b$ estimada por contagem empírica quando o estado codificado aparece ao menos duas vezes no conjunto avaliado (72,5\% das transições de teste do \emph{split} fixo) e pela distribuição \emph{softmax} do BC nos demais; $\pi_b$ é truncada em 0,01. O estimador duplamente robusto (DR), também auto-normalizado, soma o termo direto $\hat{r}(\mathbf{s}, \pi(\mathbf{s}))$ e a correção por importância com o resíduo $r - \hat{r}(\mathbf{s}, a)$, em que $\hat{r}$ é um regressor de recompensa imediata treinado na partição de treino (Tabela~\ref{tab:hiperparametros_apendice}). Os intervalos de confiança de 95\% provêm de \emph{bootstrap} de 1.000 reamostras por transição. O tamanho efetivo de amostra (ESS) das razões de importância fica entre cerca de 230 e 315 para todas as políticas, exceto o PPO (cerca de 50). Uma variante do DR com a função $Q$ do CQL como modelo direto foi calculada, mas está na escala de retorno descontado ($\gamma = 0{,}95$) e não é comparável às demais colunas; não é reportada no texto.

\section*{Configuração Base e Hiperparâmetros Executados}

Os quatro algoritmos e o regressor $\hat{r}$ usam um MLP de duas camadas ocultas, $d_{\mathcal{S}} \rightarrow h \rightarrow h \rightarrow 11$, com $h = 256$ no DQN e $h = 128$ nos demais, otimizador Adam, lote de 32 e parada antecipada com $\delta_{\min} = 10^{-4}$. A especificação do MDP prevê 17 ações; os experimentos usaram o catálogo de 11 tarefas. A dimensão de entrada é $d_{\mathcal{S}} = 34$ na comparação entre algoritmos e $d_{\mathcal{S}} = 37$ no experimento da componente cultural, com o \textit{one-hot} de país. A recompensa é a refinada do código, idêntica para todos (Equação~\ref{eq:recompensa_refinada}), sem multiplicador da Missão~6. A Tabela~\ref{tab:hiperparametros_apendice} registra os valores efetivamente executados, gerados a partir dos arquivos de configuração da execução. O orçamento de ajuste não foi equalizado entre os algoritmos: o DQN teve a grade de $\lambda_{\text{reg}}$; os demais usaram a configuração padrão dos programas.

\begin{table}[htbp]
\centering
\caption[Hiperparâmetros das políticas e do regressor]{Hiperparâmetros executados em cada política e no regressor $\hat{r}$ (iguais em toda a grade de $\lambda_{\text{reg}}$ no DQN).}
\label{tab:hiperparametros_apendice}
\small\setlength{\tabcolsep}{3.5pt}
\begin{tabular}{lccccc}
\hline
\textbf{Parâmetro} & \textbf{DQN} & \textbf{PPO} & \textbf{CQL} & \textbf{BC} & $\hat{r}$ \\
\hline
Neurônios por camada oculta ($h$) & 256 & 128 & 128 & 128 & 128 \\
Taxa de aprendizado & $3 \times 10^{-3}$ & $2{,}5 \times 10^{-4}$ & $10^{-3}$ & $10^{-3}$ & $10^{-3}$ \\
Tamanho do lote & 32 & 32 & 32 & 32 & 32 \\
Épocas máximas & 80 & 40 & 60 & 40 & 40 \\
Paciência (parada antecipada) & 12 & 8 & 8 & 8 & 8 \\
$\gamma$ & 0,9 & 0,95 & 0,95 & --- & --- \\
\textit{Dropout} & 0,1 & --- & 0,0 & --- & --- \\
$\lambda_{\text{reg}}$ & grade; 0,5 adotado & --- & --- & --- & --- \\
$\epsilon_{\text{clip}}$ & --- & 0,2 & --- & --- & --- \\
Coeficiente do crítico ($c_1$) & --- & 0,5 & --- & --- & --- \\
Coeficiente de entropia ($c_2$) & --- & 0,001 & --- & --- & --- \\
$\alpha_{\text{CQL}}$ & --- & --- & 0,2 & --- & --- \\
\textit{Weight decay} & --- & --- & 0,0 & --- & --- \\
Atualização da rede alvo & 5 épocas & --- & 20 épocas & --- & --- \\
Perda & Huber $+\ \lambda_{\text{reg}}\,$CE & recorte PPO & MSE & CE & MSE \\
Critério de seleção & Top-1 val.\ & Top-1 val.\ & perda val.\ & CE val.\ & MSE val.\ \\
Épocas treinadas (média) & 29 & 12 & 32 & 39 & 20 \\
\hline
\end{tabular}

\begin{minipage}{0.95\textwidth}
{\footnotesize Fonte: elaborado pelo autor, a partir dos arquivos de configuração dos experimentos sobre o conjunto sintético original. Nota: traço indica parâmetro inexistente na formulação ou não registrado pelo programa. ``Épocas treinadas'' é a média entre sementes no \textit{split} fixo (DQN em $\lambda_{\text{reg}} = 0{,}5$); as demais políticas registram a semente 42. Épocas treinadas pelo DQN nos demais pontos da grade: 15 ($\lambda = 0$), 21 (0,01), 15 (0,025), 17 (0,05), 21 (0,075), 24 (0,1), 26 (0,25), 29 (1), 30 (2), 29 (5) e 29 (10).}
\end{minipage}
\end{table}

O PPO, tal como implementado, é imitação ponderada por vantagem com recorte: a razão $\pi/\pi_{\text{old}}$ é avaliada apenas nas ações registradas, a vantagem é $r + \gamma V(\mathbf{s}_{t+1}) - V(\mathbf{s}_t)$ com $V$ ajustada nos mesmos dados (crítico com tronco compartilhado), e não há correção \textit{off-policy} em relação à política de comportamento. Ele colapsa (Top-1 de $0{,}234 \pm 0{,}036$) e é reportado por completude.

\section*{Componentes da Recompensa Refinada}

Em todos os experimentos sobre o conjunto sintético original, a recompensa efetivamente calculada (Equação~\ref{eq:recompensa_refinada}) é
\[
R(\mathbf{s}_t, a_t) = \mathrm{clip}_{[0,1]}\bigl(0{,}45\, r_{\text{base}} + 0{,}25\, r_{\text{nov}} + 0{,}20\, r_{\text{prog}} + 0{,}10\, r_{\text{miss}} + m_{M2}(a_t)\bigr),
\]
com os seguintes componentes:

\begin{itemize}
    \item $r_{\text{base}}$: valor de recompensa gravado pelo gerador sintético em cada transição. O \textit{pipeline} não o calcula a partir da recompensa composta; como o gerador não foi preservado, a sua composição é desconhecida
    \item $r_{\text{nov}}$: vale 1 quando o passo conclui uma tarefa ainda não concluída, e 0 nos demais casos
    \item $r_{\text{prog}} = \min(\Delta_{\text{pontos}}/100,\ 1)$, com $\Delta_{\text{pontos}} \geq 0$
    \item $r_{\text{miss}}$: vale 1 quando a missão da ação é concluída no passo (3,1\% das transições com o codificador corrigido). É evento distinto do registrado por $r_{\text{concl}}$ na recompensa composta, que é a conclusão da tarefa recomendada no passo $t$
    \item $m_{M2}(a_t) = +0{,}05$ para as ações da Missão~2, e 0 para as demais
\end{itemize}

\noindent A recompensa composta da Equação~\ref{eq:recompensa}, com $\lambda_1 = 0{,}35$ ($r_{\text{concl}}$, conclusão), $\lambda_2 = 0{,}30$ ($r_{\text{eng}}$, engajamento), $\lambda_3 = 0{,}20$ ($r_{\text{aprend}}$, aprendizagem) e $\lambda_4 = 0{,}15$ ($r_{\text{cult}}$, adequação cultural), é a especificação para a operação com aprendizes. Ela não foi computada nos experimentos \textit{offline}. O coeficiente $\lambda_{\text{reg}}$ da regularização por comportamento do DQN (Equação~\ref{eq:dqn_reg_comportamento}) não deve ser confundido com esses pesos.

\section*{Seleção de $\lambda_{\text{reg}}$}

A grade foi $\{0;\ 0{,}01;\ 0{,}025;\ 0{,}05;\ 0{,}075;\ 0{,}1;\ 0{,}25;\ 0{,}5;\ 1;\ 2;\ 5;\ 10\}$, com cinco sementes por ponto no \emph{split} fixo e a época de parada escolhida na validação. A seleção é feita na validação, por dois critérios. O valor de Top-1 máximo de validação, com desempate pelo menor $\lambda_{\text{reg}}$, é $\lambda^{*} = 5$ (Top-1 de validação $0{,}879 \pm 0{,}004$). O ponto de operação recomendado é $\lambda_{\text{1EP}} = 0{,}5$, o menor $\lambda_{\text{reg}}$ cujo Top-1 de validação fica a um erro-padrão binomial, $\sqrt{p(1-p)/n_{\text{val}}} = 0{,}016$, do máximo; o Top-1 de validação satura a partir desse ponto (0,871 em 0,5; 0,879 em 5). Com o erro-padrão entre sementes (0,0016) no lugar do binomial, a regra seleciona o próprio $\lambda^{*}$, porque as sementes quase não alteram o Top-1 de validação. Em $\lambda_{\text{1EP}} = 0{,}5$, o teste do \emph{split} fixo dá Top-1 de $0{,}825 \pm 0{,}006$, Top-3 de $0{,}994 \pm 0{,}002$, recompensa associada de $0{,}573 \pm 0{,}002$, SNIS de $0{,}542 \pm 0{,}007$ e DR de $0{,}554 \pm 0{,}003$. Na validação cruzada, a seleção aninhada produz $\lambda_{\text{1EP}}$ de 0,25, 0,5, 0,25, 0,25 e 0,25 e $\lambda^{*}$ de 2, 0,5, 1, 1 e 1 nos cinco \emph{folds}. A grade completa consta da Tabela~\ref{tab:q_lambda}, e a leitura da fronteira fidelidade--recompensa, da Seção~\ref{sec:q_regularizacao}.

\section*{Configuração do Experimento da Componente Cultural}
\label{sec:apendice_cultural}

O experimento usa o codificador corrigido, o \emph{split} fixo por perfil e o DQN regularizado em $\lambda^{*} = 5$, com cinco sementes. As duas configurações partilham tudo, inclusive o \textit{one-hot} de país no estado ($d_{\mathcal{S}} = 37$), e diferem apenas na recompensa:

\begin{itemize}
    \item \textbf{Recompensa}: $R$ da seção anterior (controle) ou a mistura convexa $R_{\text{cult}} = 0{,}85\, R + 0{,}15\, f_{\text{cult}}(a_t, \text{país})$. O termo cultural não entra como parcela $\lambda_4 r_{\text{cult}}$ da recompensa composta
    \item \textbf{Função de adequação cultural}: $f_{\text{cult}}$ da Equação~\ref{eq:f_cult}, instanciada conforme a Tabela~\ref{tab:fcult}, com valor neutro 0,5 quando o país do aprendiz não está registrado. Como o Brasil vale 1,0 e Portugal vale 0,5 em todas as missões, a função só discrimina ações no par (Angola, M6): a mistura altera a ordenação de ações apenas para aprendizes de Angola, e só entre a M6 e as demais missões
    \item \textbf{Atribuição de país aos perfis}: alternada entre Brasil, Angola e Portugal sobre os 25 perfis (9, 8 e 8), sem correlação com literacia, segmento, sentimento ou qualquer outro atributo. No teste do \emph{split} fixo há 150 transições de Angola, 150 do Brasil e 75 de Portugal. Por isso o conjunto não permite decidir a questão da eficácia (Seção~\ref{sec:q_cultural})
    \item \textbf{Teste de divergência}: para cada transição de teste, o \textit{one-hot} de país é substituído pelas três atribuições possíveis, mantido idêntico o restante do vetor, e o estado conta como divergente quando ao menos duas atribuições produzem ações distintas
    \item \textbf{Resultados}: Top-1 de teste de $0{,}810 \pm 0{,}008$ sem $f_{\text{cult}}$ e $0{,}812 \pm 0{,}009$ com; divergência por país de $12{,}3 \pm 2{,}3\%$ sem e $15{,}0 \pm 2{,}5\%$ com. Teste $t$ de Welch entre as configurações: $t = 1{,}83$, $p = 0{,}106$; pareado por semente, $t = 2{,}10$, $p = 0{,}103$ ($n = 5$ por grupo). A Tabela~\ref{tab:q_divergencia} reúne os valores
\end{itemize}

\section*{Reprodutibilidade}
\label{sec:apendice_reprodutibilidade}

\begin{itemize}
    \item \textbf{Execução}: um comando, \texttt{python run\_all.py}, na pasta \path{drl/experimentos_2026-10-10/}, regenera \texttt{resultados.json}, o relatório \texttt{RESULTADOS.md} e as três figuras da Seção~\ref{sec:q_drl}. Tempo total de 7,4 minutos em CPU (\emph{split} fixo 3,0 min; validação cruzada 3,8 min; experimento cultural 0,6 min), em Python 3.12.10 com torch 2.14.1 (CPU) e numpy 2.5.3.
    \item \textbf{Sementes}: 42 a 46 no \emph{split} fixo e 42 na validação cruzada, fixadas para o gerador de números pseudoaleatórios do \textit{framework}, para a inicialização dos pesos e para o embaralhamento dos lotes. Os resultados do \emph{split} fixo são média e desvio padrão entre sementes; os da validação cruzada, média e desvio padrão entre \emph{folds}, com intervalos por \emph{bootstrap} das 2.000 transições fora da amostra.
    \item \textbf{Protocolo de seleção e reporte}: a época de parada e o valor de $\lambda_{\text{reg}}$ são escolhidos na validação; as métricas provêm do teste (\emph{split} fixo) ou das transições fora da amostra (validação cruzada). A partição de treino é a única sobre a qual há atualização de pesos.
    \item \textbf{Artefatos e depósito}: o conjunto preservado compreende os programas de treinamento, de seleção de $\lambda_{\text{reg}}$, de avaliação \textit{off-policy} e do experimento cultural, os arquivos de configuração, o conjunto de transições, as partições por perfil e um arquivo de métricas por execução. O que será depositado, o que é irrecuperável e o prazo do depósito constam da Seção~\ref{subsec:disponibilidade}. O programa gerador do conjunto sintético não foi preservado.
    \item \textbf{Configuração do sistema nos pilotos}: modelo de linguagem Llama~4 Maverick (\texttt{meta-llama/Llama-4-Maverick-17B-128E-Instruct}), com pesos quantizados em 4 bits e execução em ONNX Runtime, acessado por API em \textit{endpoint} próprio na nuvem da equipe; temperatura de 0,1, até 3.500 \textit{tokens} por resposta e contexto de 32.768 \textit{tokens}; orquestração em LangGraph 1.1.10. A geração aumentada por recuperação usa os \textit{embeddings} \texttt{cnmoro/bert-tiny-embeddings-english-portuguese}, trechos de 256 \textit{tokens} com sobreposição de 25, PostgreSQL~18 com pgvector~0.8.2 (índice IVFFlat, similaridade de cosseno), filtro por país e $k = 5$ trechos recuperados. Nos pilotos, o conteúdo das conversas não foi enviado a provedor externo de modelo de linguagem; trafegou pelos servidores do WhatsApp e pela infraestrutura de nuvem da equipe.
\end{itemize}

\noindent Esses elementos tornam os experimentos reproduzíveis, mas não alteram o caráter sintético dos dados (Seção~\ref{sec:q_drl}).

\section*{Nota histórica: a primeira execução}

A primeira execução dos experimentos de DRL, de agosto de 2026, tinha quatro problemas. O codificador atribuía ao estado seguinte, em todas as transições, literacia 0, segmento ``iniciante'' e indicadores de missão zerados, o que corrompia o alvo de Bellman de DQN, CQL e PPO e mantinha $r_{\text{miss}}$ sempre em 0. A partição era por interação (1.581, 206 e 213 transições), com 66,2\% dos estados de teste presentes no treino. O valor de $\lambda_{\text{reg}}$ ($0{,}05$), por sua vez, fora escolhido pelo Top-3 observado na partição de teste. Além disso, o multiplicador de recompensa da Missão~6 ($1{,}5$) estava implementado apenas nos programas de DQN e PPO, de modo que a comparação entre algoritmos dependia de nova execução sob recompensa uniforme. Os números então reportados (DQN com Top-1 de 0,574 e recompensa de 0,709; BC 0,901; CQL 0,809; $\lambda_{\text{reg}} = 0{,}05$ com 0,832, 1,000 e 0,563; divergência por país de 8,7\% contra 7,2\%) não são comparáveis aos atuais e não devem ser lidos como resultado. Naquele mesmo teste, a regra fixa ``próxima tarefa do catálogo'' já obtinha Top-1 de 0,592 e recompensa de 0,708, acima do DQN, o que motivou a revisão do protocolo. A regularização por comportamento, a análise de teto e o teste de divergência por país vêm daquela execução; a regra sequencial, a partição por perfil, a avaliação \textit{off-policy} e a seleção de $\lambda_{\text{reg}}$ na validação foram introduzidas com a correção. Os programas e dados da primeira execução estão preservados sem alteração em pasta própria.

\section*{Simulador Cognitivo}
\label{sec:apendice_simulador}

O simulador cognitivo da Seção~\ref{sec:q_simulador} complementa o conjunto sintético original como ambiente de avaliação. Nele, cada aprendiz simulado tem conhecimento latente, e o desfecho é o ganho de Hake num pós-teste simulado, distinto da recompensa de treino.

\noindent\textbf{Modelo.} O catálogo é o de 11 tarefas dos experimentos (M1, 5; M2, 2; M3, 2; M6, 2), com uma habilidade por tarefa e, como pré-requisito, a tarefa anterior na cadeia. As 11 tarefas são sempre escolhíveis; os pré-requisitos agem pelo modelo cognitivo, e não por bloqueio. A maestria de cada habilidade segue um BKT estendido com maestria contínua, pré-requisitos, esquecimento e abandono. O parâmetro $h \in \{0;\ 0{,}5;\ 1\}$ controla a heterogeneidade entre segmentos e aprendizes: com $h = 0$, todos os aprendizes têm os parâmetros médios. A Tabela~\ref{tab:apendice_simulador} resume os parâmetros. O episódio termina após $H = 15$ decisões, no abandono ou quando todas as habilidades atingem maestria de 0,95. O currículo fixo escolhe a primeira tarefa não concluída na ordem do catálogo e, com todas concluídas, revisa as tarefas 1 a 4.

\begin{table}[htbp]
\centering
\caption[Parâmetros do simulador cognitivo]{Parâmetros e distribuições do simulador cognitivo ($h$: heterogeneidade; dp: desvio padrão).}
\label{tab:apendice_simulador}
\footnotesize\setlength{\tabcolsep}{4pt}
\begin{tabular}{lp{0.68\textwidth}}
\hline
\textbf{Componente} & \textbf{Especificação} \\
\hline
Segmentos & iniciante, intermediário e avançado, nas proporções 9:11:5 dos perfis do conjunto sintético \\
$P(L_0)$ & médias por segmento 0,15, 0,35 e 0,60, contraídas para a média geral $\bar{\mu}$ por $\bar{\mu} + h(\mu_{\text{seg}} - \bar{\mu})$; aprendiz $\sim$ Beta (dp $0{,}1h$); habilidade $\sim$ Beta em torno do aprendiz (dp $0{,}2h$) \\
$P(T)$ & $c_T \times (0{,}75;\ 1{,}0;\ 1{,}25)$ por segmento, com $c_T = 0{,}6586$ (calibrado); dp $0{,}1h$ entre aprendizes e $0{,}05h$ entre habilidades \\
$P(S)$, $P(G)$, $P(F)$ & Beta com médias 0,1, 0,2 e 0,01 por passo; dp $0{,}04h$, $0{,}06h$ e $0{,}005h$ entre aprendizes \\
Pré-requisitos & $P(T)$ efetivo multiplicado por 0,3 se a maestria média dos pré-requisitos for menor que 0,6 \\
Esquecimento & $L \leftarrow L - P(F)(L - L_0)$ por passo sem exposição \\
Abandono & probabilidade por passo $0{,}005 + \alpha(\mathit{Fr} + \mathit{Bo})$, com $\alpha = 0{,}0108$ (calibrado); $\mathit{Fr}$ e $\mathit{Bo}$ acumulam frustração (pré-requisito não dominado) e tédio (maestria acima de 0,9), com decaimento 0,5 \\
Sentimento observável & ansioso se $\mathit{Fr} \geq 0{,}75$; cansado se $\mathit{Bo} \geq 0{,}75$; motivado se acertou a última tarefa; neutro nos demais casos; 25\% de ruído uniforme \\
Desfecho & pré e pós-teste = média de $P(\text{acerto}) = L(1-S) + (1-L)G$ nas 11 habilidades; pós-teste no passo $H$ para todos, inclusive quem abandona (intenção de tratar); retenção 30 passos depois; Hake $g = (\text{pós} - \text{pré})/(1 - \text{pré})$ \\
\hline
\end{tabular}

{\footnotesize Fonte: elaborado pelo autor.}
\end{table}

\noindent\textbf{Calibração.} A única âncora empírica é o Hake do tutor no piloto brasileiro (0,51). O parâmetro $c_T$ foi ajustado para que o currículo tenha Hake de 0,500 em $h = 0{,}5$, e $\alpha$, para que o abandono do currículo seja de 15\% em $h = 1$; este segundo alvo é escolha de desenho, e não dado. Os dois ajustes usam bisseções alternadas numa população própria (semente 999, 20.000 aprendizes). Na calibração, o currículo tem Hake de 0,542, 0,500 e 0,436 e abandono de 7,1\%, 9,8\% e 15,0\% em $h = 0$, $0{,}5$ e 1. Os demais parâmetros são plausíveis, e não estimados. A sensibilidade a três deles (multiplicador de pré-requisito em $\{0{,}2;\ 0{,}3;\ 0{,}5\}$, limiar de pré-requisito em $\{0{,}5;\ 0{,}6;\ 0{,}7\}$ e escala de todos os dp em $\{0{,}75;\ 1;\ 1{,}25\}$) foi testada em 27 combinações em $h = 1$, com $c_T$ recalibrado em cada uma. A sensibilidade a $\alpha$, $P(S)$, $P(G)$, $P(F)$ e às médias de $P(L_0)$ não foi testada.

\noindent\textbf{Observação e recompensa.} O agente não observa a maestria verdadeira. O vetor de observação tem 46 dimensões: as 34 da representação dos experimentos de DRL, a maestria estimada de cada uma das 11 habilidades e o passo normalizado. A estimativa vem de um rastreador BKT com maestria contínua (crença Beta por habilidade), que usa os parâmetros médios do segmento e a mesma estrutura do simulador, o que favorece as políticas adaptativas, e recebe um pré-teste observado de 2 itens por habilidade. A recompensa de treino é o aprendizado estimado pelo rastreador no passo, menos 0,01 por passo e menos 1,0 no abandono. Nenhuma política é treinada nem selecionada pelo Hake.

\noindent\textbf{Políticas.} Foram comparados: o currículo fixo; uma política aleatória que respeita os pré-requisitos concluídos; LinUCB ($\alpha = 1$) e Thompson linear ($v = 0{,}5$), treinados em 2.000 aprendizes e depois congelados; o DQN \textit{offline} treinado com registros do currículo com 20\% de ações uniformes (cerca de 28 mil transições de treino e 7 mil de validação), com a configuração do DQN deste apêndice, exceto lote de 256, 40 épocas e paciência 8, e $\lambda_{\text{reg}}$ escolhido na grade $\{0;\ 0{,}01;\ 0{,}05;\ 0{,}1;\ 0{,}25;\ 0{,}5;\ 1;\ 5\}$ pelo Top-1 com regra de um erro-padrão ou pela SNIS de validação; o DQN \textit{online} (Double DQN, MLP $256 \times 2$, taxa de aprendizado $10^{-3}$, $\gamma = 0{,}9$, lote 256, memória de 200.000 transições, 24.000 aprendizes de treino, $\epsilon$ de 1,0 a 0,05 na primeira metade, \textit{checkpoint} escolhido pelo retorno de treino numa população de validação); a escolha gulosa pelo rastreador, como referência implementável; e o oráculo míope, que vê a maestria verdadeira, como teto. Cada política enfrenta, em cada nível de $h$, cinco populações de teste de 1.000 aprendizes, disjuntas das de treino e validação, com números aleatórios comuns, o que permite diferenças pareadas. Os intervalos de 95\% vêm de \textit{bootstrap} de aprendizes (2.000 reamostras). A validação da avaliação \textit{off-policy} usa registros do currículo com 20\% de ações uniformes e $\pi_b$ conhecida, e compara SNIS e DR com o valor verdadeiro de cada política no simulador.

\noindent\textbf{Planejador LLM.} O planejador avaliado no simulador é o Claude Opus 5.5 (Anthropic), executado como subagente que decide por rodadas. A cada rodada, um programa grava em arquivo a descrição textual do estado de cada aprendiz ativo, e o subagente devolve um arquivo com uma tarefa (letra de A a K) por aprendiz. O texto contém apenas o que o sistema em operação veria: segmento, sentimento declarado, passo, acertos no pré-teste, e, por tarefa, o pré-requisito, o número de vezes que foi feita e o domínio estimado pelo rastreador em três níveis (baixo, abaixo de 0,4; médio, até 0,7; alto); seguem as três últimas respostas. O subagente não tem acesso ao simulador, ao seu código nem à maestria verdadeira, e decide cada aprendiz apenas pelo texto do seu estado. Uma letra ausente ou inválida aplica a tarefa do currículo e é contada como inválida. A avaliação usa 200 aprendizes (semente 42) em $h = 0{,}5$ e em $h = 1$; o currículo é aplicado aos mesmos aprendizes, com os mesmos números aleatórios, e a diferença de Hake é pareada. O planejador do sistema em operação é o Llama~4 Maverick; o Claude Opus 5.5 é usado aqui como limite de capacidade de um planejador LLM de fronteira, sem pretensão de equivalência entre os dois.

\noindent\textbf{Reprodutibilidade.} Na pasta \path{drl/experimentos_2026-10-11_simulador/}: \texttt{python run\_all.py} regenera a calibração, a comparação de políticas, a validação da avaliação \textit{off-policy}, o relatório e as figuras (17,5 minutos em CPU, com Python 3.12.10, numpy 2.5.3 e torch 2.14.1); \texttt{python sensibilidade.py}, executado depois, refaz as 27 combinações (16,9 minutos). O planejador LLM usa \texttt{python planejador\_rodadas.py init}, um passo \texttt{step} por rodada de decisões e \texttt{finalize}, que calcula o contraste pareado com o currículo; a visão em texto de cada rodada é gerada por \texttt{rodada\_para\_texto.py}.

%%%%%%%%%%%%%%%%%%%%%%%%%%%%%%%%%%%%%%%%%%%%%%%%%%%%%%%%%%%%%%%%%%%%%%%%%%%%%%%%
